\section{Related Work}
\label{sec:related}

Four adjacent lines of work inform the design of \eps: empirical
studies of deprecated API usage (which motivate the problem), token-%
level and sequence-level calibration of LLMs for code (which share
the underlying logprob signal), multi-sample uncertainty
quantification (which is the dominant alternative and the closest
point of comparison), and adaptive decoding (which uses the same
per-token entropy signal for a different purpose). We review each in
turn, closing with static analysis.

\subsection{Deprecated and version-split API usage}
\label{subsec:related-deprecated}

Wang et al.~\cite{wang2025deprecated} present the first systematic
study of deprecated API usage in LLM-based code completion, evaluating
seven models on 145 API mappings drawn from eight popular Python
libraries across $28{,}125$ completion prompts. They measure API
Usage Plausibility (AUP) and Deprecated Usage Rate (DUR), finding
that deprecated-API emission is prevalent across all evaluated
models; the root cause they identify is that deprecated API usages
are present in training data while deprecation status itself is not
reliably propagated to inference time. Their mitigations (REPLACEAPI
and INSERTPROMPT) operate at the prompt or post-generation level and
require a known deprecation mapping.

Their finding establishes that the problem we address is common in
practice rather than a corner case: on version-split APIs and renamed
SDK methods, models routinely commit to a deprecated branch without
surface indication. Our work is complementary. Wang et al.\ measure
\emph{how often} LLMs emit deprecated code and diagnose \emph{why}
the output looks confident; \eps measures the \emph{signal at
generation time} that distinguishes a confidently-chosen branch from
an internally-contested one, and --- crucially for the review loop
--- localizes that signal to a specific token rather than producing
only a function-level score.

\subsection{Token-level and sequence-level calibration}
\label{subsec:related-calibration}

Spiess et al.~\cite{spiess2025calibration} study calibration and
correctness of language models for code, computing several confidence
signals (including the sequence-level mean token probability, which
we refer to throughout as \papg) and evaluating whether they
correlate with functional correctness on HumanEval and related
benchmarks. They find that generative code models are not well-%
calibrated out of the box and that Platt scaling can recover
calibration in settings where labelled correctness data is available.
The \papg signal --- the exponential of the mean log-probability over
all generated tokens --- is derived from the same raw logprob data
that \eps uses, and is the closest signal-level comparator for the
function-level part of \eps's output.

Four differences between that setup and ours deserve emphasis. First,
\emph{aggregation}: \papg takes a sequence-level mean of per-token
probabilities; \eps takes a max after AST filtering, so a single
high-entropy API-selector token is not diluted by long confident
stretches. Second, \emph{filtering}: \eps restricts the input set to
semantically meaningful tokens (attribute accesses, keyword arguments,
import targets), while \papg averages over every token including
whitespace and cosmetic names. Third, \emph{prompt regime}: Spiess et
al.\ evaluate on HumanEval and similar benchmarks where correctness
is well-defined by unit tests; we evaluate on version-split API
prompts where the ground truth is ambiguous enough to require LLM
judging (\S\ref{subsec:gt-disclosure}). Fourth, \emph{output
granularity}: \papg produces one scalar per function and is silent on
\emph{which token} carried the uncertainty. This fourth distinction
is the one that matters for the review loop (\S\ref{sec:comparison}).
Because \papg and \eps operate on the same underlying logprobs at the
same $1\times$ API cost, a direct head-to-head comparison is
possible; we report it in \S\ref{sec:comparison}.

Two concurrent code-specific studies further establish that token-level
uncertainty localizes to the same regions \eps targets. Sharma and
David~\cite{sharma2025multisample} adapt entropy and mutual information
to code generation and show, using symbolic execution as ground truth,
that uncertainty-based abstention reduces incorrect outputs to near
zero; their finding grounds the core premise of \eps in prior empirical
evidence without requiring symbolic execution at inference time. Gros
and Devanbu~\cite{gros2025localized} assign calibrated uncertainty to
individual code regions using a supervisor model, demonstrating that
sub-function localization is both feasible and informative. Our
contribution relative to both is that we infer token-level risk
directly from the generator's logprob distribution---no auxiliary model
or test executor required---and route the result to a downstream review
loop rather than to abstention.

Earlier token-level calibration studies (Kadavath et
al.~\cite{kadavath2022language}; Varshney et
al.~\cite{varshney2023stitch}) establish that LLMs are reasonably
calibrated at the token level for factual questions and that
selective generation can exploit this. Gupta et
al.~\cite{gupta2024cascades} prove that sequence-level aggregation of
token probabilities suffers from length bias---longer outputs are
systematically over- or under-penalised by mean logprob---and show
that token-level deferral rules substantially outperform na\"{i}ve
sequence-level means; this provides a principled explanation for why
\papg is a challenging baseline to beat despite sharing the same raw
logprob signal. Santilli et al.~\cite{santilli2025length} further
demonstrate that length bias compromises \emph{evaluation} of
uncertainty methods across eight signals and four models, validating
our choice to compare \eps and \papg at matched operating points
rather than by raw AUC. Our work extends that line to
code, where AST filtering is what makes per-token entropy actionable
rather than dominated by cosmetic variance. Conformal prediction
approaches~\cite{quach2023conformal,angelopoulos2022conformal}
provide set-valued predictions with coverage guarantees but do not
localize which token within a generation is uncertain; TECP~\cite{xu2025tecp}
applies split conformal prediction directly to token-entropy scores,
providing per-output coverage guarantees that could in future work be
used to calibrate the \eps threshold with statistical guarantees. Uncertainty
propagation through multi-stage ML pipelines is studied in adjacent
domains~\cite{jungo2020analyzing}.

\subsection{Multi-sample and behavioural-diversity methods}
\label{subsec:related-multisample}

Inter-generational methods sample the same prompt multiple times and
measure disagreement among outputs. Kuhn et
al.~\cite{kuhn2023semantic} introduce semantic entropy: multiple
completions are clustered by semantic equivalence (via natural
language inference), and cross-cluster entropy provides a
hallucination signal for factual QA.
UQLM~\cite{uqlm2024} packages semantic negentropy and related signals
as an open-source library; we include it here as the closest off-%
the-shelf tool and report its performance in \S\ref{sec:comparison}
in one sentence. Beyond these, a growing line of
work~\cite{sharma2025multisample} applies stronger notions of
equivalence --- up to and including symbolic execution or semantic-%
equivalence checks on generated code --- to improve behavioural
consistency detection specifically for code outputs.

These methods are structurally limited for the task we study:
intra-generational uncertainty does not surface as output variation
when the token distribution is peaked, as it typically is for code.
A model can be internally torn between two API calls and emit one of
them unanimously across re-rolls; we provide the empirical instance
in \S\ref{subsec:multisample-comp}. A stronger equivalence check
improves precision on Type 3 (implementation-approach) false
positives, where two semantically equivalent implementations would
be conflated by output-level diversity alone; but it cannot recover
a signal that never varies across samples in the first place. The
intra/inter distinction is not a claim that inter-generational
methods are badly implemented, only that they are measuring a
different quantity.

\subsection{Generation-time intervention and adaptive decoding}
\label{subsec:related-adadec}

AdaDec~\cite{adadec2025} reads the same per-token Shannon entropy
that \eps reads, but uses it at decoding time rather than review
time: when the entropy at the current step exceeds a learned
threshold, the decoder pauses and triggers a lookahead re-ranking
step over candidate continuations. On HumanEval and MBPP across
model sizes from $0.6$B to $8$B, AdaDec reports up to $15.5$pp
improvement in Pass@1 over greedy decoding while pausing on
$3.95$--$11.88\%$ of tokens --- the same order of magnitude as
\eps's flagged-token rate of $10.5\%$ (\S\ref{subsec:token-focus}),
which is not a coincidence: both identify the same set of decision
points. The signal is the same as ours; the action is different.
AdaDec improves the model's own output at generation time; \eps
surfaces residual uncertainty \emph{after} the model has committed,
for routing to downstream review. The two approaches are compatible:
an AdaDec-enhanced generator still produces a token distribution at
each step, and \eps operates on whatever distribution the final
decoder emits. We do not evaluate a combined pipeline but expect
\eps values to be lower on average after AdaDec-style intervention,
with residual high-\eps tokens concentrating on consequential
decisions that even a lookahead re-rank could not resolve.
Consequently, applying AdaDec upstream of our pipeline is expected to
reduce average \eps values but not eliminate high-entropy API-selector
tokens, which remain the primary target of the review loop.

\subsection{Entropy-guided refinement and self-repair}
\label{subsec:related-refinement}

A separate line of work uses uncertainty signals to trigger iterative
repair rather than review. Self-Refine~\cite{madaan2023selfrefine}
establishes the canonical pattern: a single model generates, critiques
its own output, and refines iteratively, achieving $\approx 20\%$
absolute improvement across diverse tasks without supervision. The
refinement loop is unconditional---every generation is sent back
regardless of confidence.

Conditional entropy-triggered variants tighten this by gating
second-pass work on measured uncertainty. Correa and de
Matos~\cite{correa2025entropyloop} compute Shannon entropy over top-$k$
alternatives at each step, trigger refinement when entropy, perplexity,
or low-confidence token counts exceed learned thresholds, and feed an
uncertainty report back to the originating model; they report
$16$\,pp accuracy gains at one-third the cost of reasoning-class
models. ERGO~\cite{khalid2025ergo} applies the same principle to
multi-turn generation, resetting the prompt context when entropy
spikes sharply across turns.

Our review loop shares the entropy-triggered gating with these methods
but differs in two ways. First, the second-pass agent is a separate
reviewer with domain-specific prompting, not a self-critique call to
the originating model; this removes the correlation between generator
failure and reviewer failure. Second, the review is over
\emph{code-specific risk}---the flagged token and its alternatives,
not the full response---which concentrates reviewer attention and
enables the precision improvement from $27\%$ to $94\%$
(\S\ref{sec:comparison}).

\subsection{Static analysis, linters, and execution-based evaluation}
\label{subsec:related-static}

Static analysis and linter rules catch \emph{known}-bad patterns for
specific libraries once deprecation has been encoded into a ruleset.
These approaches are complementary to \eps and represent the floor
of detection: when a rule exists, a linter will catch the problem
whether or not the model hesitated. \eps operates in the gap --- it
fires where the model itself was uncertain, which includes version
splits not yet in any ruleset (a library updated last week, a half-%
migrated codebase) and internal renames that static analysis has no
fact to test against. Execution-based code-generation
benchmarks~\cite{chen2021evaluating,austin2021program,%
jimenez2024swebench} measure whether the output passes test cases,
not whether the model was confident; a silent dependency update can
turn a fully-passing generation into a regression that no test in
the current suite exercises. To our knowledge, no prior work treats
within-generation \emph{token-level} uncertainty as a first-class
routing signal for code review that localizes where within a
function to look; the closest is CodeT~\cite{chen2022codet} and
related self-consistency approaches, which are inter-generational by
construction, and post-generation repair with uncertainty-based
ranking~\cite{ye2022selfapr}, which operates on already-generated
candidates.
